\documentclass{article}
\title{Riemannian metrics}
\author{Math 6702}
\date{Week 1}

\begin{document}
\maketitle

We start by fixing the basic objects of the course: Riemannian metrics and the geometric quantities they induce, so that length, angle, and volume have precise meanings on a manifold.

\section*{Reading}
From \textit{Riemannian Geometry} by Sylvestre Gallot, Dominique Hulin, and Jacques Lafontaine, read \textsection II.A Existence Theorems and First Examples: ``Definitions'' and ``First examples.''

\end{document}
